\par\Needspace{4\baselineskip}
\begin{center}
\SourceLeaf{024b}\pdfbookmark[2]{6.}{article-006}6.
\end{center}
Wir wenden uns nun dem Zusammenhang zwischen der Bewegung und der \emph{Zeit} zu. Setzen wir wie in Art. 1 den in der Zeit $t$ um die Sonne überstrichenen Flächeninhalt $= \frac{1}{2}g$ und die Masse des bewegten Körpers $= \mu$, wobei die Masse der Sonne $= 1$ gesetzt wird, so erhalten wir $g = kt\BareRoot{p}.\BareRoot{(1+\mu)}$. Das Differential des Flächeninhalts ist $= \frac{1}{2}rr\,\mathrm{d}v$, woraus $kt\BareRoot{p}.\BareRoot{(1+\mu)}
= \int rr\,\mathrm{d}v$ folgt. Das Integral ist so zu nehmen, dass es für $t=0$ verschwindet. Die Integration muss je nach Art des Kegelschnitts unterschiedlich durchgeführt werden. Daher betrachten wir die einzelnen Arten getrennt und beginnen mit der ELLIPSE.

Der Bruch, durch den $r$ aus $v$ bestimmt wird, hat zwei Glieder im Nenner. Diesen Nachteil beseitigen wir zunächst, indem wir anstelle von $v$ eine neue Größe einführen. Dazu setzen wir $\operatorname{tang}\frac{1}{2}v\BareRoot{\frac{1-e}{1+e}}
= \operatorname{tang}\frac{1}{2}E$. Die letzte Formel des vorigen Artikels ergibt dann für $r$
\[
r = \frac{p\cos\frac{1}{2}E^2}{(1+e)\cos\frac{1}{2}v^2} = p\left(\frac{\cos\frac{1}{2}E^2}{1+e}+\frac{\sin\frac{1}{2}E^2}{1-e}\right) = \frac{p}{1-ee}(1-e\cos E).
\]
Weiter erhalten wir $\frac{\mathrm{d}E}{\cos\frac{1}{2}E^2} = \frac{\mathrm{d}v}{\cos\frac{1}{2}v^2}\BareRoot{\frac{1-e}{1+e}}$ und damit $\mathrm{d}v = \frac{p\,\mathrm{d}E}{r\BareRoot{(1-ee)}}$. Daraus folgt
\[
rr\,\mathrm{d}v = \frac{r\,\grave{p}\,\mathrm{d}E}{\BareRoot{(1-ee)}} = \frac{pp}{(1-ee)^{\frac{3}{2}}}(1-e\cos E)\,\mathrm{d}E,
\]
und durch Integration
\[
kt\BareRoot{p}.\BareRoot{(1+\mu)} = \frac{pp}{(1-ee)^{\frac{3}{2}}}(E-e\sin E)+\mathrm{Const.}
\]
Rechnen wir also die Zeit vom Durchgang durch das Perihel an, wo $v=0$ und $E=0$ gelten
\SourceLeaf{025}und somit die Konstante $= 0$ ist, so erhalten wir wegen $\frac{p}{1-ee}=a$
\[
E-e\sin E = \frac{kt\BareRoot{(1+\mu)}}{a^{\frac{3}{2}}}\cdot
\]

In dieser Gleichung muss der Hilfswinkel $E$, der \emph{exzentrische Anomalie} heißt, im Bogenmaß ausgedrückt werden. Natürlich können wir den Winkel weiterhin in Grad usw. angeben, sofern wir auch $e\sin E$ und $\frac{kt\BareRoot{(1+\mu)}}{a^{\frac{3}{2}}}$ auf dieselbe Weise ausdrücken. Durch Multiplikation mit $206264{,}806$ erhält man diese Größen in Bogensekunden. Die Multiplikation der zweiten Größe können wir uns ersparen, indem wir $k$ von vornherein in Bogensekunden einsetzen und somit anstelle des oben angegebenen Werts $k = 3548^{\prime\prime}{,}18761$ setzen. Sein Logarithmus ist $= 3{,}550\,0065\,746$. — Die auf diese Weise ausgedrückte Größe $\frac{kt\BareRoot{(1+\mu)}}{a^{\frac{3}{2}}}$ heißt \emph{mittlere Anomalie}. Sie wächst also proportional zur Zeit, und zwar täglich um $\frac{k\BareRoot{(1+\mu)}}{a^{\frac{3}{2}}}$. Dieser Zuwachs heißt \emph{mittlere tägliche Bewegung}. Die mittlere Anomalie bezeichnen wir mit $M$.
